\bookchapter{Приложение E. Благодарности и пояснения}

Спасибо, что дочитали этот «телефонный курс» до конца! При написании мы опирались на множество официальных инструкций и руководств производителей, а также общедоступных научно-популярных технических статей, чтобы обеспечить точность информации.

Технологии меняются каждый день: содержание книги опирается на популярные технологии и продукты на момент написания. Когда вы будете её читать, возможно, появятся новые функции и устройства. Надеемся, что идеи и методы, которым учит эта книга, помогут вам спокойно встречать будущие цифровые новинки.

\booksection{Основные обновления версии 1.6}

\begin{itemize}
  \item Расширена тема «Ценность цифровых активов и принципы управления», добавлена глава «Управление фотографиями и файловой системой» (именование «время + событие», разбор по категориям, автоочистка через 72 часа, обезличивание EXIF).
  \item Доработана система защиты «чистый/защищённый режим» и защита от «цифровых паразитов».
  \item Добавлены распространённые в эпоху ИИ схемы мошенничества (подмена лица и голоса / выдача себя за авторитет / фиктивные инвестиции) и подходы к защите от них.
  \item Переписаны разделы о защите от доксинга («вскрытия») и атак социальной инженерии.
  \item Расширена глава о доступности (полное управление VoiceOver/TalkBack, увеличение и цветовые фильтры, описания изображений с ИИ, субтитры в реальном времени, двухканальные звонки RTT/TTY).
  \item Рамка промптов переписана по схеме «роль — контекст — задача — ограничения», с 3 полными примерами сценариев.
  \item Добавлены три технических приложения продвинутого уровня: Дополнение к главе 1: разбор расширения памяти на Android-телефонах (определение по сути, принципы zRAM/swap, узкое место по скорости, границы применимости, названия у разных производителей и способы включения).
  \item Дополнение к главе 20: Termux — переносимая Linux-терминальная среда на Android-устройствах (ключевые функции, экосистема плагинов, установка и настройка, типичные сценарии: мобильная разработка / администрирование серверов / тестирование безопасности).
  \item Дополнение к главе 23: вводный гид по GitHub — от хостинга кода к платформе совместной работы (разница Git и GitHub, ключевые понятия, регистрация и первые шаги, поток сотрудничества через Pull Request, напоминания о безопасности, способы доступа из Китая).
\end{itemize}

\booksection{Основные обновления версии 1.7}

1. Добавлены 4 новые главы

\begin{itemize}
  \item Глава 29: цифровое здоровье — анализ экранного времени, лимиты приложений, режим «Не беспокоить», правило 20-20-20 для глаз, спасение шеи, управление сном, «избирательная слепота» при информационной тревожности, недельный челлендж цифрового здоровья.
  \item Глава 30: трансляция экрана и автомобиль — сравнение трёх лагерей AirPlay / Miracast / DLNA, варианты проводной трансляции, десктопный режим Samsung DeX, четыре автомобильных решения CarPlay / Android Auto / HiCar / CarLife, практика трансляции в переговорной.
  \item Глава 31: создание коротких видео — три вопроса к теме, 5 шагов настройки съёмки, пять приёмов движения камеры (наезд, отъезд, панорама, сдвиг, сопровождение), шестиэтапный рабочий процесс CapCut, ключевые моменты обложки, контент-принцип «первые 3 секунды решают всё».
  \item Глава 32: справочник цифровой безопасности в экстренных случаях — SOP «золотых 30 минут» после потери телефона, порядок действий за 72 часа после обмана, проверка в безопасном режиме при заражении, шесть шагов возврата украденного аккаунта, реакция на сообщения-шантаж, пять ежедневных мер профилактики.
\end{itemize}

2. Добавлены две новые главы

\begin{itemize}
  \item Глава 27: забота о старших и родных — сравнительная таблица упрощённых режимов по брендам, шесть главных схем мошенничества и три линии обороны, удалённая помощь (ChangLian / Sunflower), режим заботы WeChat и голосовой ввод, SOS / детекция падения / напоминания о лекарствах.
  \item Глава 28: гид по выбору телефона — метод трёх шагов, таблица параметров «смотрите только пять», рекомендации по пяти бюджетным уровням, таблица выбора по семи сценариям, пять главных ловушек при покупке, анализ времени покупки и каналов.
\end{itemize}

3. Новые функции озвучки и взаимодействия

\begin{itemize}
  \item Плавающая панель управления озвучкой: воспроизведение/пауза, предыдущий/следующий фрагмент, регулировка скорости 0.5x–2.5x, остановка.
  \item Быстрая кнопка наушников справа: озвучка запускается одним нажатием и автоматически начинается с текущего места чтения.
  \item Горячие клавиши: R — включить/выключить озвучку, ESC — остановить.
  \item Улучшение озвучки: автоматический непрерывный переход по главам, элементы управления озвучкой автоматически скрываются при печати.
\end{itemize}

4. Обновление структуры

Оглавление, учебный путь (четыре шага пять), карта глав, номер версии (1.61.7), нумерация приложений (пять шесть) и метаданные — всё синхронно обновлено.

5. Обновление актуальности содержимого

\begin{itemize}
  \item Стандарты процессоров: Snapdragon 8 Elite Gen5 / Dimensity 9500 отнесены к топовому классу (базис 2026 года).
  \item Референсные флагманы: iPhone 17 Pro, Samsung S26 Ultra, Xiaomi 17 Ultra, Huawei Pura 80 Ultra, vivo X300 Pro.
  \item Главы об ИИ: обновлены до DeepSeek R1, ChatGPT-5 и стандарта производительности встроенных NPU в 80 TOPS.
  \item Складные экраны: ценовой диапазон опустился до уровня 4000–5000 юаней.
  \item Гид по выбору: все бюджетные уровни и сценарные рекомендации обновлены до рынка 2026 года.
  \item Сеть и экосистема: Wi-Fi 7 / 5G-A / экосистема HarmonyOS / LE Audio — всё обновлено до последних стандартов.
\end{itemize}

6. Новый функциональный модуль: учебная панель (расположена под блоком Hero)

\begin{itemize}
  \item Кольцевая диаграмма прогресса чтения: в реальном времени показывает процент прочитанных глав, с анимацией SVG-кольца.
  \item Самопроверка здоровья телефона: чек-лист из 8 пунктов безопасности (пароль, двухфакторная аутентификация, поиск телефона, облачная резервная копия, резервная копия WeChat, безопасность платежей, приложение против мошенничества, отключение платежей без пароля).
  \item Сохранение состояния: отметки записываются в localStorage, оценка здоровья вычисляется динамически и отображается градиентной полосой прогресса.
  \item Подсказки по уровням: в зависимости от балла показываются разные отзывы вроде «максимум / неплохо / сдал / риски безопасности».
\end{itemize}

\booksection{Основные обновления версии 1.8}

Добавлена часть 7 «Практика программирования с ИИ» (главы 37–47, всего 11 глав): революция парадигм и Agentic Engineering, пять принципов мышления в ИИ-программировании, контекст-инжиниринг и разработка на основе спецификаций (SDD), инженерия требований, панорама инструментов ИИ-программирования 2026 года, панорама больших языковых моделей и практика API, диалоговый инжиниринг, полный путь проекта от 0 до 1, обеспечение качества, путь вглубь (MCP / взаимодействие нескольких агентов) и справочник. Синхронно обновлены оглавление, учебный путь (пять шагов шесть), карта глав и быстрый индекс. Материал основан на книге «Мышление ИИ-программирования: методологическое руководство от нуля до full-stack практики», издание августа 2026 года.

Добавлены часть 8 «Практика разработки под Android» (главы 48–58, всего 11 глав) и часть 9 «Практика разработки под iOS» (главы 59–69, всего 11 глав): обе платформы охватывают настройку окружения (Android Studio / Xcode), экспресс-курс языка (Kotlin / Swift), фреймворки интерфейса (Jetpack Compose / SwiftUI), данные и сеть, два полноценных практических проекта (книга учёта + запрос погоды), тестирование и публикацию (Google Play / App Store) и справочник. Синхронно обновлены оглавление, учебный путь (шесть шагов восемь), карта глав и быстрый индекс. Материал основан на книгах «Разработка под Android с нуля» и «Разработка под iOS с нуля», издания августа 2026 года.

\booksection{Дополнения к содержанию: полезные программы и блокировка назойливых звонков}

\begin{itemize}
  \item В главу 13 добавлен шестой раздел «Полезные системные инструменты»: инструменты наведения порядка и очистки вроде Qingzhuo, встроенная очистка хранилища, Национальный центр против мошенничества, фирменные «менеджеры телефона», защита операторов от назойливых звонков и AdGuard, а также принцип выбора «сначала системные / дополнительно официальные / сторонние — сдержанно».
  \item Сравнительная таблица встроенной блокировки назойливых звонков по брендам: охватывает пути входа и рекомендуемые к включению пункты для iOS, Xiaomi/Redmi, Huawei/Honor, OPPO/realme/OnePlus, vivo/iQOO и Samsung.
  \item Расширены возможности блокировки назойливых звонков в главе 8: добавлены практические инструкции по распознаванию звонящих, чёрным и белым спискам, фильтрации SMS по ключевым словам, заглушению неизвестных звонков в iOS, связке с режимом «Не беспокоить» и другим, с перекрёстными ссылками на главу 13.
  \item Трёхуровневая защита: системная блокировка + защита оператора от назойливых звонков + Национальный центр против мошенничества, с акцентом на белый список, чтобы не блокировать нужное.
\end{itemize}

\booksection{Дополнения к содержанию: погодные предупреждения / предупреждения о стихийных бедствиях / предупреждения о землетрясениях}

\begin{itemize}
  \item В главу 13 добавлен седьмой раздел: погодные предупреждения (официальные метеоканалы + системное погодное приложение + цветовые уровни), предупреждения о стихийных бедствиях (множество каналов: метеорология / чрезвычайные службы), принцип и заблуждения о предупреждениях о землетрясениях (предупреждение $\neq$ прогноз).
  \item Таблица включения предупреждений о землетрясениях по брендам: типичные пути входа и нюансы разрешений для Xiaomi, Huawei/Honor, устройств OPPO, vivo, iPhone, Samsung и других.
  \item В главу 32 добавлен «сценарий ноль»: порядок действий после сигнала о землетрясении, реакция на ливень/тайфун, связь со спутниковым SOS / связью без сети; справочник экстренных действий перекрёстно ссылается на настройки предупреждений в главе 13.
  \item Связка с погодным разделом: в пункте о погодных приложениях в теме поездок добавлены предупреждения-уведомления и напоминание «уведомления должны быть включены, иначе вы их не получите».
\end{itemize}

\booksection{Дополнения: связка устройств и инструменты продуктивности}

\begin{itemize}
  \item Раздел 8 главы 13: управление компьютером со смартфона (Sunflower, ToDesk) и главные правила защиты от мошенничества; ИИ для работы (Trae / WorkBuddy / Dumate) и сценарий управления ПК со смартфона; живые обои HoYoDial от miHoYo.
  \item Быстрая передача по локальной сети и буфер обмена: LocalSend, FeiGe Chuanshu, AirDrop; KDE Connect, универсальный буфер обмена, обмен между фирменными телефонами и компьютерами.
  \item Одноразовые коды 2FA: Aegis, Microsoft Authenticator, резервные коды и нюансы при смене телефона.
  \item TickTick и сканирование бумажных документов: подробнее о ведении расписания и задач; дополнены инструменты сканирования и приёмы фотосъёмки.
\end{itemize}

Желаем вам в совершенстве освоить смартфон, наслаждаться удобством и радостями технологий и стать осознанным и уверенным цифровым гражданином.

— Ваш проводник в цифровую жизнь
